% https://github.com/sb2nov/resume (MIT-licensed template format).
% Resume content: Jegannarayanan Sivakasiulaganathan.
% Compile with: pdflatex -interaction=nonstopmode Jegannarayanan_Sivakasiulaganathan_Compiler_Resume.tex

\documentclass[letterpaper,11pt]{article}

\usepackage[T1]{fontenc}
\usepackage[utf8]{inputenc}
\usepackage{lmodern}
\usepackage[empty]{fullpage}
\usepackage{titlesec}
\usepackage{enumitem}
\usepackage[hidelinks]{hyperref}
\usepackage{fancyhdr}
\usepackage[english]{babel}
\usepackage{tabularx}
\input{glyphtounicode}
\pdfgentounicode=1

\pagestyle{fancy}
\fancyhf{}
\renewcommand{\headrulewidth}{0pt}
\renewcommand{\footrulewidth}{0pt}
\setlength{\footskip}{10pt}

% Single-column, compact layout adapted from sb2nov/resume.
\addtolength{\oddsidemargin}{-0.47in}
\addtolength{\evensidemargin}{-0.47in}
\addtolength{\textwidth}{0.94in}
\addtolength{\topmargin}{-0.43in}
\addtolength{\textheight}{0.86in}
\setlength{\tabcolsep}{0pt}
\setlength{\parindent}{0pt}
\raggedright
\urlstyle{same}

\titleformat{\section}{\vspace{-4pt}\scshape\raggedright\large}{}{0em}{}[\titlerule\vspace{-5pt}]
\titlespacing*{\section}{0pt}{11pt}{6pt}

\newcommand{\resumeSubheading}[4]{%
  \begin{tabular*}{\textwidth}[t]{@{}l@{\extracolsep{\fill}}r@{}}
    \textbf{#1} & \small #2 \\
    \textit{\small #3} & \textit{\small #4} \\
  \end{tabular*}\vspace{-2pt}
}
\newcommand{\resumeProjectHeading}[2]{%
  \begin{tabular*}{\textwidth}[t]{@{}l@{\extracolsep{\fill}}r@{}}
    \textbf{#1} & \small #2 \\
  \end{tabular*}\vspace{-3pt}
}
\newenvironment{resumeBullets}{%
  \begin{itemize}[leftmargin=0.18in, itemsep=3pt, topsep=3pt, parsep=0pt, partopsep=0pt]
  \small
}{\end{itemize}\vspace{2pt}}

\begin{document}

\begin{center}
	{\LARGE\bfseries Jegannarayanan Sivakasiulaganathan}\\[3pt]
	\small Dedham, MA \enspace $\vert$ \enspace
	\href{tel:+16692331124}{+1-669-233-1124} \enspace $\vert$ \enspace
	\href{mailto:Jegannarayanan_Sivakasiulaganathan@student.uml.edu}{Jegannarayanan\_Sivakasiulaganathan@student.uml.edu}\\[2pt]
	\small \href{https://github.com/handel-fan}{github.com/handel-fan}
\end{center}
\vspace{-9pt}

\section{Education}

\resumeSubheading
{University of Massachusetts Lowell}{Expected May 2027}
{M.S. Computer Science}{}

\vspace{5pt}

{\small
	\textbf{Relevant Coursework:}
	Compiler Engineering, Effective Functional Programming,
	Graduate Algorithms, Operating Systems I,
	Deep Learning (LLMs/LVMs).
}

\vspace{12pt}

\resumeSubheading
{University of Massachusetts Amherst}{May 2018}
{B.S. Computational and Applied Mathematics}{}

\section{Research and Compiler Projects}
\resumeProjectHeading{Master's Thesis: Lean TensorLib to MLIR Lowering \textnormal{\emph{(in progress)}}}{Lean 4, MLIR $\mid$ 2026--Present}
\begin{resumeBullets}
	\item Investigating a compiler path from high-level operations in Lean's TensorLib to
	MLIR's \texttt{linalg} dialect, with matrix multiplication as an initial case
	study.
	\item Examining Lean's elaborated \texttt{Expr} representation to identify TensorLib
	calls, arguments, and types before generating lower-level IR.
	\item Wrote corresponding Lean and MLIR matrix-multiplication examples; examining
	semantic-preservation requirements for tensor operations and lowering.
\end{resumeBullets}

\resumeProjectHeading{LLVM/Clang AST Exploration}{C++, LLVM, Clang}
\begin{resumeBullets}
	\item Built LLVM/Clang from source and implemented Clang AST traversal tooling using
	\texttt{RecursiveASTVisitor}, \texttt{ASTConsumer}, and \texttt{FrontendAction}
	to inspect C++ declarations.
\end{resumeBullets}

\resumeProjectHeading{Monkey Programming Language Interpreter}{C++ $\mid$ \href{https://github.com/handel-fan/first-interpreter-in-cpp}{GitHub}}
\begin{resumeBullets}
	\item Implementing a Monkey-language interpreter in C++ to explore lexical analysis,
	parsing, and interpreter architecture.
\end{resumeBullets}

\section{Professional Experience}
\resumeSubheading{KBRA}{Dresher, PA (Remote)}{Software Developer (promoted from Associate Software Developer)}{Aug 2022--Present}
\begin{resumeBullets}
	\item Led technical-debt initiatives and migration away from an unsupported PHP
	framework toward C\#/.NET services, improving maintainability and reducing
	reliance on legacy code.
	\item Developed C\# ASP.NET Core backend services, worked with Hot Chocolate GraphQL
	schemas, and contributed to React frontend features.
	\item Tested Docker startup-time retrieval of AWS Secrets Manager configuration for a
	.NET service, without persisting secrets to disk.
\end{resumeBullets}

\resumeSubheading{Aspen Technology}{Bedford, MA}{Software Developer (previously Software Developer I)}{May 2018--Aug 2022}
\begin{resumeBullets}
	\item Maintained and developed C\# desktop software, resolving defects and supporting
	multiple production applications.
	\item Progressed from Software Developer I to Software Developer in July 2021.
\end{resumeBullets}

\section{Technical Skills}
 {\small
  \textbf{Programming:} Lean 4, C++, Python, C\#, Haskell, TypeScript; Go (exploratory)\\[2pt]
  \textbf{Compilers and formal methods:} MLIR, LLVM/Clang, AST traversal, Rocq (coursework)\\[2pt]
  \textbf{Software and systems:} Linux, Git, Docker, GitLab CI/CD, ASP.NET Core, React
 }

\end{document}
